\documentclass{article}
\usepackage[T1]{fontenc}
\usepackage[utf8]{inputenc}
\usepackage[french]{babel}
\usepackage{longtable}
\usepackage{booktabs}
\usepackage{listings}
\usepackage{xcolor}
\usepackage{float}
\usepackage{graphicx}

\lstset{
  basicstyle=\ttfamily\footnotesize,
  breaklines=true,
  breakatwhitespace=true,
  columns=fullflexible,
  keepspaces=true,
  frame=single,
  numbers=left,
  numberstyle=\tiny,
  captionpos=b,
  showstringspaces=false
}

\title{Analyse des thèmes dans les romans de Zola}
\author{Morgan RICHARD}

\date{}

\begin{document}

\maketitle

\section{Introduction}

Émile Zola, l'un des écrivains les plus célèbres du XIXe siècle, est connu pour ses romans naturalistes qui explorent les aspects sociaux, économiques et psychologiques de la vie humaine. 

Dans cette analyse, nous allons examiner les thèmes récurrents dans les œuvres de Zola et déterminer s'ils apparaissent dans plusieurs de ses romans. Des méthodes d'analyse textuelle seront utilisées pour identifier les thèmes principaux à travers les différents romans.

Ce travail est le premier d'une série d'analyses portant sur un corpus de textes littéraires issus de différents auteurs naturalistes et non naturalistes.

Deux méthodes d'analyse ont été utilisées pour identifier les thèmes dans les romans de Zola : 
\vspace{0.5cm}
\begin{itemize}
\item Le NMF (Non-negative Matrix Factorization)\footnote{N. Gillis, \textit{Nonnegative Matrix Factorization}, Society for Industrial and Applied Mathematics, coll. « Data Science », 2020.} : une technique de réduction de dimensionnalité qui permet d'extraire des thèmes à partir d'un corpus de textes. 

\item K-means : une méthode de clustering qui permet de regrouper les thèmes récurrents en fonction de leur similarité thématique.
\end{itemize}
\vspace{0.5cm}
Ces méthodes seront expliquées en détail dans les sections suivantes, où nous présenterons les résultats de l'analyse thématique des romans de Zola.

\newpage
\section{Présentation du corpus}

Comme indiqué en introduction, ce travail porte sur un corpus de textes de Zola. Le choix a été fait de n'inclure que ses romans dans cette analyse, afin de se concentrer sur les thèmes récurrents dans l'œuvre de cet auteur. Aucune nouvelle n'a été utilisée pour cette analyse, car les romans de Zola sont plus longs et offrent une richesse thématique plus importante que les nouvelles. De plus, un déséquilibre entre ces deux types de textes pourrait fausser les résultats de l'analyse thématique.

Le corpus comprend les romans suivants :

\begin{longtable}{p{1.6cm}p{4.2cm}p{8cm}}
\caption{Corpus des œuvres de Zola utilisées dans l'analyse}\\
\toprule
\textbf{Date} & \textbf{Titre} & \textbf{Fichier} \\
\midrule
\endfirsthead

\toprule
\textbf{Date} & \textbf{Titre} & \textbf{Fichier} \\
\midrule
\endhead

1865 & La confession de Claude & \texttt{1865\_La\_confession\_de\_Claude.\_clean.txt} \\
1866 & Le vœu d'une morte & \texttt{1866\_Le\_voeu\_d\_une\_morte.\_clean.txt} \\
1867 & Les mystères de Marseille & \texttt{1867\_Les\_mysteres\_de\_Marseille.\_clean.txt} \\
1867 & Thérèse Raquin & \texttt{1867\_Therese\_Raquin.\_clean.txt} \\
1868 & Madeleine Férat & \texttt{1868\_Madeleine\_Ferat.\_clean.txt} \\
1871 & La fortune des Rougon & \texttt{1871\_1\_La\_fortune\_des\_Rougon.\_clean.txt} \\
1872 & La curée & \texttt{1871\_2\_La\_curee.\_clean.txt} \\
1873 & Le ventre de Paris & \texttt{1873\_3\_Le\_ventre\_de\_Paris.\_clean.txt} \\
1874 & La conquête de Plassans & \texttt{1874\_4\_La\_conquete\_de\_Plassans.\_clean.txt} \\
1875 & La faute de l'abbé Mouret & \texttt{1875\_5\_La\_faute\_de\_l\_abbe\_Mouret.\_clean.txt} \\
1876 & Son Excellence Eugène Rougon & \texttt{1876\_6\_Son\_Excellence\_Eugene\_Rougon.\_clean.txt} \\
1877 & L'assommoir & \texttt{1877\_7\_L\_assommoir.\_clean.txt} \\
1878 & Une page d'amour & \texttt{1878\_8\_Une\_page\_d\_amour.\_clean.txt} \\
1880 & Nana & \texttt{1880\_9\_Nana.\_clean.txt} \\
1882 & Pot-Bouille & \texttt{1882\_10\_Pot-bouille.\_clean.txt} \\
1883 & Au Bonheur des dames & \texttt{1883\_11\_Au\_Bonheur\_des\_dames.\_clean.txt} \\
1884 & La joie de vivre & \texttt{1884\_12\_La\_joie\_de\_vivre.\_clean.txt} \\
1885 & Germinal & \texttt{1885\_13\_Germinal.\_clean.txt} \\
1886 & L'œuvre & \texttt{1886\_14\_L\_oeuvre.\_clean.txt} \\
1887 & La terre & \texttt{1887\_15\_La\_terre.\_clean.txt} \\
1888 & Le rêve & \texttt{1888\_16\_Le\_reve.\_clean.txt} \\
1890 & La bête humaine & \texttt{1890\_17\_La\_bete\_humaine.\_clean.txt} \\
1891 & L'argent & \texttt{1891\_18\_L\_argent.\_clean.txt} \\
1892 & La débâcle & \texttt{1892\_19\_La\_debacle.\_clean.txt} \\
1893 & Le docteur Pascal & \texttt{1893\_20\_Le\_docteur\_Pascal.\_clean.txt} \\
1894 & Lourdes & \texttt{1894\_1\_Lourdes.\_clean.txt} \\
1896 & Rome & \texttt{1896\_2\_Rome.\_clean.txt} \\
1898 & Paris & \texttt{1898\_3\_Paris.\_clean.txt} \\
1899 & Fécondité & \texttt{1899\_1\_Fecondite.\_clean.txt} \\
1901 & Travail & \texttt{1901\_2\_Travail.\_clean.txt} \\
1903 & Vérité & \texttt{1903\_3\_Verite.\_clean.txt} \\

\bottomrule
\end{longtable}

Nous constatons que le corpus est composé de 31 romans, correspondant aux principaux romans publiés de Zola :
\vspace{0.3cm}
\begin{itemize}
    \item 5 : romans antérieurs aux \textit{Rougon-Macquart}
    \item  20 :  Rougon-Macquart
    \item  3 : Trois Villes
    \item  3 : Quatre Évangiles publiés
    
\end{itemize}
\vspace{0.3cm}
Pour ce qui est des \textit{Quatre Évangiles}, \textit{Vérité} paraît à titre posthume. \textit{Justice} ne paraîtra jamais, l'ouvrage étant resté à l'état d'ébauche au moment de la mort de l'écrivain. 
\newpage
\section{Méthodologie de préparation des données}

\subsection{Préparation des données brutes en données propres}

Les romans au format \texttt{.txt} correspondent à des versions nettoyées, c'est-à-dire à des textes dont les éléments non textuels ont été supprimés. Ces éléments comprennent notamment les titres, les chapitres, les notes de bas de page et d'autres informations éditoriales qui pourraient perturber l'analyse. Ce nettoyage permet d'obtenir une représentation plus précise des thèmes présents dans les romans.

De plus, le choix a été fait de conserver une seule phrase par ligne. Ce format facilite l'analyse textuelle, notamment pour les méthodes de NMF et de K-means, qui nécessitent une structure de données homogène pour fonctionner correctement.

Cette logique de formatage a également été appliquée en prévision des analyses futures. Les romans issus de prochains traitements pourront ainsi être intégrés de manière cohérente au corpus. Ce choix permet de limiter les différences de structure entre les fichiers et d'obtenir un format homogène pour l'analyse sans trop complexifier le processus de préparation des données.

\begin{longtable}{p{8cm}r}
\caption{Nombre de tokens par roman dans l'ordre chronologique}  \\
\toprule
\textbf{Titre du roman} & \textbf{Nombre de tokens} \\
\midrule
\endfirsthead

\toprule
\textbf{Titre du roman} & \textbf{Nombre de tokens} \\
\midrule
\endhead

La confession de Claude & 47253 \\
Le vœu d'une morte & 38722 \\
Les mystères de Marseille & 125253 \\
Thérèse Raquin & 66363 \\
Madeleine Férat & 97318 \\
La fortune des Rougon & 116685 \\
La curée & 105120 \\
Le ventre de Paris & 110442 \\
La conquête de Plassans & 113245 \\
La faute de l'abbé Mouret & 114667 \\
Son Excellence Eugène Rougon & 126186 \\
L'assommoir & 158340 \\
Une page d'amour & 102185 \\
Nana & 141573 \\
Pot-Bouille & 134844 \\
Au Bonheur des dames & 147521 \\
La joie de vivre & 117255 \\
Germinal & 164473 \\
L'œuvre & 130613 \\
La terre & 163687 \\
Le rêve & 66379 \\
La bête humaine & 125418 \\
L'argent & 143010 \\
La débâcle & 187046 \\
Le docteur Pascal & 112713 \\
Lourdes & 173511 \\
Rome & 233958 \\
Paris & 177389 \\
Fécondité & 220131 \\
Travail & 198712 \\
Vérité & 223274 \\

\bottomrule
\end{longtable}

Le tableau fournit une description quantitative du corpus en indiquant, pour chaque roman, le nombre de tokens obtenus après nettoyage. Les œuvres sont présentées dans l'ordre chronologique, ce qui permet d'observer l'évolution du volume textuel au sein de la production romanesque de Zola. Les écarts sont importants entre les romans courts du début de carrière, comme \textit{Le vœu d'une morte}, et les œuvres plus volumineuses de la fin de carrière, notamment \textit{Rome}, \textit{Fécondité}, \textit{Travail} et \textit{Vérité}.

Cette information est importante pour l'interprétation des résultats, car la longueur des textes peut avoir un effet sur la représentation des thèmes. Un roman plus long fournit mécaniquement davantage d'occurrences lexicales et davantage de segments analysables. La segmentation en paquets de phrases vise donc à produire des unités d'analyse plus comparables, tout en conservant suffisamment de contexte pour identifier les thèmes dominants.
\subsection{Préparation du corpus pour l'analyse thématique}

Pour effectuer l'analyse sur le corpus des romans de Zola, nous avons procédé de la façon qui semblait la plus pertinente pour obtenir des résultats fiables et significatifs. 

Les romans ont été segmentés en paquets de phrases. Les paquets étaient composés de 40 phrases, ce qui correspond à une taille de texte suffisante pour capturer les thèmes présents dans les romans, tout en permettant une analyse thématique plus fine. Bien entendu, ce choix de taille de paquet est arbitraire et pourrait être ajusté en fonction des besoins de l'analyse. Plusieurs tailles de paquets ont été testées\footnote{Différents essais ont été réalisés avec des paquets de 10, 50, 70 et 100 phrases.} et les paquets de 40 phrases ont été retenus, car ils permettent d'obtenir des résultats plus cohérents et significatifs que les autres tailles testées. Ce choix relève d'un arbitrage méthodologique fondé sur des essais empiriques et sur une analyse qualitative des résultats obtenus avec différentes tailles de paquets.


\vspace{0.3cm}
Le code suivant permet de segmenter les romans en paquets de 40 phrases.
\vspace{0.3cm}
\begin{lstlisting}[language=Python, caption={Code de segmentation du texte en paquets de phrases}]
def segmenter_en_paquets(texte, taille_paquet=40):
    # Segmente le texte en paquets de phrases
    phrases = texte.splitlines()
    phrases = [p.strip() for p in phrases if p.strip()]
    paquets = []

    for i in range(0, len(phrases), taille_paquet):
        paquet = " ".join(phrases[i:i+taille_paquet])
        paquets.append(paquet)
    return paquets

rows = []

for _, row in df.iterrows():
    # Iteration sur chaque ligne du DataFrame
    paquets = segmenter_en_paquets(row["texte_brut"])

    for idx, paquet in enumerate(paquets):
        rows.append({
            "nom_fichier": row["nom_fichier"],
            "id_paquet": idx,
            "phrases_paquet": paquet
        })

df_phrases = pd.DataFrame(rows)
\end{lstlisting}

\subsection{Statistiques descriptives du corpus}

\begin{longtable}{p{8cm}r}
\caption{Nombre de paquets de phrases par roman} \\
\toprule
\textbf{Titre du roman} & \textbf{Nombre de paquets} \\
\midrule
\endfirsthead

\toprule
\textbf{Titre du roman} & \textbf{Nombre de paquets} \\
\midrule
\endhead

La confession de Claude & 65 \\
Le vœu d'une morte & 60 \\
Les mystères de Marseille & 184 \\
Thérèse Raquin & 80 \\
Madeleine Férat & 124 \\
La fortune des Rougon & 154 \\
La curée & 123 \\
Le ventre de Paris & 136 \\
La conquête de Plassans & 165 \\
La faute de l'abbé Mouret & 171 \\
Son Excellence Eugène Rougon & 191 \\
L'assommoir & 228 \\
Une page d'amour & 166 \\
Nana & 225 \\
Pot-Bouille & 207 \\
Au Bonheur des dames & 198 \\
La joie de vivre & 164 \\
Germinal & 225 \\
L'œuvre & 163 \\
La terre & 222 \\
Le rêve & 82 \\
La bête humaine & 154 \\
L'argent & 155 \\
La débâcle & 221 \\
Le docteur Pascal & 135 \\
Lourdes & 198 \\
Rome & 231 \\
Paris & 206 \\
Fécondité & 259 \\
Travail & 208 \\
Vérité & 227 \\

\bottomrule
\end{longtable}

Ce tableau indique le nombre de paquets de phrases obtenus pour chaque roman après segmentation du corpus. Les écarts observés reflètent principalement les différences de longueur entre les œuvres : les textes les plus courts, comme \textit{Le vœu d'une morte} ou \textit{La confession de Claude}, produisent peu de paquets, tandis que les romans plus volumineux, comme \textit{Fécondité}, \textit{Rome} ou \textit{L'assommoir}, en produisent davantage.

\subsubsection{Analyse des statistiques descriptives du nombre de paquets par roman}
\begin{table}[H]
\centering
\caption{Statistiques descriptives du nombre de paquets par roman}
\begin{tabular}{lr}
\toprule
\textbf{Indicateur} & \textbf{Valeur} \\
\midrule
Nombre de romans & 31 \\
Moyenne & 171.84 \\
Écart-type & 52.40 \\
Minimum & 60 \\
Premier quartile & 145 \\
Médiane & 171 \\
Troisième quartile & 214.5 \\
Maximum & 259 \\
\bottomrule
\end{tabular}
\end{table}


Les statistiques descriptives du nombre de paquets par roman permettent d'évaluer la structure quantitative du corpus après segmentation. Le corpus comprend 31 romans, avec une moyenne d'environ 172 paquets de phrases par œuvre. La médiane, située à 171 paquets, est très proche de cette moyenne. Cette proximité indique que la distribution générale du nombre de paquets n'est pas fortement asymétrique : la majorité des romans se situe autour d'un volume intermédiaire, même si certains textes s'écartent nettement de cette tendance.

L'écart-type, d'environ 52 paquets, montre toutefois une dispersion importante autour de la moyenne. Cette variation s'explique principalement par les différences de longueur entre les romans. Les œuvres les plus courtes, comme \textit{Le vœu d'une morte}, produisent seulement 60 paquets, tandis que les textes les plus longs, comme \textit{Fécondité}, atteignent 259 paquets. L'écart entre le minimum et le maximum est donc notable, puisqu'il existe presque un rapport de un à quatre entre le roman le moins segmenté et le roman le plus segmenté.

Les quartiles permettent de préciser cette répartition. Le premier quartile, situé à 145 paquets, indique qu'un quart des romans contient 145 paquets ou moins. Le troisième quartile, situé à 214,5 paquets, indique qu'environ trois quarts des romans restent sous ce seuil. La moitié centrale du corpus se situe donc entre 145 et 214,5 paquets, ce qui montre que la majorité des œuvres conserve une taille relativement comparable après segmentation. Les valeurs extrêmes ne remettent donc pas en cause l'homogénéité générale du corpus, mais elles doivent être prises en compte dans l'interprétation des résultats.

Le coefficient de variation est d'environ 30 \%, puisqu'il correspond au rapport entre l'écart-type et la moyenne. Cette valeur signale une variabilité modérée : les romans ne sont pas tous de taille équivalente, mais les écarts restent attendus pour un corpus littéraire composé d'œuvres publiées à des périodes différentes et appartenant à plusieurs cycles romanesques. Cette dispersion est donc méthodologiquement acceptable, à condition de garder à l'esprit que les romans les plus longs fournissent davantage d'unités d'analyse.

Cette information est importante pour l'analyse thématique. Chaque paquet de phrases constitue une unité soumise aux méthodes de topic modeling. Par conséquent, un roman plus long contribue mécaniquement à un plus grand nombre d'observations qu'un roman court. Les thèmes présents dans les romans volumineux peuvent donc être davantage représentés dans les résultats globaux. La segmentation en paquets de 40 phrases permet néanmoins de réduire cet effet, car elle transforme les romans en unités textuelles comparables, tout en conservant suffisamment de contexte pour identifier des thèmes cohérents.

Ainsi, ces statistiques montrent que le corpus présente un équilibre satisfaisant pour une analyse thématique. Les différences de taille entre les romans existent, mais elles sont documentées et peuvent être intégrées à l'interprétation des résultats. Cette étape descriptive permet donc de mieux contrôler les effets liés à la longueur des œuvres avant l'application des méthodes de NMF et de K-means.

\subsubsection{Analyse des statistiques descriptives du nombre de mots par paquet}

\begin{table}[H]
\centering
\caption{Statistiques descriptives du nombre de mots par paquet}
\begin{tabular}{lr}
\toprule
\textbf{Indicateur} & \textbf{Valeur} \\
\midrule
Nombre de paquets & 5327 \\
Moyenne & 785.30 \\
Écart-type & 212.13 \\
Minimum & 42 \\
Premier quartile & 636 \\
Médiane & 753 \\
Troisième quartile & 902.5 \\
Maximum & 2944 \\
\bottomrule
\end{tabular}
\end{table}

Les statistiques descriptives du nombre de mots par paquet permettent d'évaluer la taille effective des unités textuelles utilisées pour l'analyse thématique. Le corpus segmenté comprend 5327 paquets, avec une moyenne d'environ 785 mots par paquet. La médiane, située à 753 mots, est relativement proche de cette moyenne, ce qui indique que la majorité des paquets se situe autour d'un volume textuel comparable.

L'écart-type est d'environ 212 mots, ce qui traduit une dispersion modérée autour de la moyenne. Le coefficient de variation est d'environ 27 \%, puisqu'il correspond au rapport entre l'écart-type et la moyenne. Cette valeur indique que les paquets ne sont pas parfaitement homogènes en longueur, mais que la variabilité reste acceptable pour une analyse thématique. Cette variation s'explique par le fait que les paquets sont construits à partir d'un nombre fixe de phrases, et non d'un nombre fixe de mots : selon le style, le rythme narratif ou la longueur syntaxique des phrases, deux paquets de 40 phrases peuvent contenir un nombre de mots assez différent.

Les quartiles permettent de préciser cette distribution. Un quart des paquets contient au maximum 636 mots, tandis que la moitié des paquets contient au maximum 753 mots. Le troisième quartile, situé à 902,5 mots, indique que 75 \% des paquets ne dépassent pas environ 903 mots. Ainsi, la majorité des unités d'analyse se concentre dans une zone relativement stable, comprise approximativement entre 636 et 903 mots pour la moitié centrale du corpus. Cette concentration est favorable au topic modeling, car elle limite les écarts de taille entre les unités comparées.

Les valeurs extrêmes doivent toutefois être prises en compte. Le paquet le plus court contient seulement 42 mots, tandis que le plus long atteint 2944 mots. Ces deux valeurs indiquent la présence d'unités atypiques. Le paquet de 42 mots correspond probablement à une fin de roman ou à un segment résiduel composé de moins de phrases que les autres. À l'inverse, le paquet de 2944 mots peut provenir d'un passage composé de phrases particulièrement longues ou d'une segmentation imparfaite du texte source. Ces cas extrêmes peuvent influencer l'analyse, car un paquet très long contient davantage d'information lexicale et peut donc peser plus fortement dans la représentation des thèmes.

Dans l'ensemble, ces statistiques montrent que la segmentation en paquets de 40 phrases produit des unités textuelles globalement comparables. Malgré une variabilité réelle, la distribution reste suffisamment concentrée autour de la médiane pour permettre une analyse thématique cohérente. Il serait toutefois pertinent d'examiner les paquets extrêmes afin de vérifier s'ils correspondent à des cas normaux du corpus ou à des problèmes de segmentation.


\subsection{Lemmatisation du corpus de Zola}

Après avoir obtenu un corpus segmenté en paquets de phrases et relativement homogène en termes de nombre de mots, la prochaine étape de préparation des données a consisté à effectuer une lemmatisation du corpus.

Pour réaliser cette lemmatisation, nous avons utilisé la bibliothèque \texttt{spaCy} en Python, qui offre des outils performants pour le traitement automatique du langage. La lemmatisation permet de regrouper les différentes formes d'un même mot sous une seule forme de base : par exemple, les mots « manger », « mange » et « manges » sont tous réduits au lemme « manger ». Cette étape facilite l'identification des thèmes et contribue à réduire la dimensionnalité du corpus.

Comme pour la segmentation en paquets de phrases, plusieurs essais ont été réalisés avec la bibliothèque de lemmatisation \texttt{spaCy}.

\begin{lstlisting}[language=Python, caption={Importation du modèle français de spaCy}]
nlp = spacy.load("fr_core_news_lg", disable=["ner", "parser"])
# Chargement du modele de langue francaise de spaCy.
# Les composants de reconnaissance des entites nommees et d'analyse
# syntaxique sont desactives afin d'accelerer le traitement.

nlp.max_length = 2_000_000
# Augmentation de la longueur maximale des textes traites par spaCy.
# Cette valeur est ajustable selon la taille des textes et la memoire
# disponible sur la machine.
\end{lstlisting}

Nous avons choisi d'utiliser le modèle \texttt{fr\_core\_news\_lg} de spaCy, qui est un modèle de langue française de grande taille. Ce choix a été motivé par la richesse lexicale et la complexité syntaxique des romans de Zola, qui nécessitent un modèle capable de gérer une grande variété de formes linguistiques. Le modèle \texttt{fr\_core\_news\_lg} a été retenu car il offre une lemmatisation plus précise que des modèles plus légers\footnote{Les modèles \texttt{fr\_core\_news\_sm} et \texttt{fr\_core\_news\_md} ont également été envisagés.}, ce qui est important pour l'identification des thèmes dans les textes littéraires.


Une liste de mots à exclure a été construite progressivement au cours de plusieurs itérations de prétraitement\footnote{Elle est présente en annexe du document.}. Elle résulte de l'observation conjointe des sorties de lemmatisation et de la matrice terme-document. L'objectif était d'identifier les formes lexicales très fréquentes dans le corpus, mais peu discriminantes pour l'analyse thématique.

Cette liste contient d'abord des mots-outils, tels que les pronoms personnels, les conjonctions, les déterminants et les prépositions. Ces termes structurent grammaticalement les phrases, mais ils ne permettent pas d'identifier les thèmes dominants des romans. Leur présence importante dans la matrice terme-document risquait donc de réduire la lisibilité des topics produits par la NMF.

La liste a également été enrichie par l'ajout de certains noms propres, notamment des noms de personnages récurrents dans les romans de Zola. Ces noms apparaissent fréquemment dans les textes, mais leur poids lexical tendait à orienter les topics vers des personnages plutôt que vers des thèmes plus généraux, comme le travail, la famille, la religion, l'argent ou les rapports sociaux. Leur exclusion permet donc de privilégier une interprétation thématique plutôt qu'une simple identification des protagonistes.

Ce filtrage relève d'un choix méthodologique empirique. Lors des premiers essais, plusieurs mots restaient fortement représentés dans les résultats malgré les paramètres appliqués au lemmatiseur et au vectoriseur. La liste d'exclusion a donc été ajustée à partir de l'examen des termes les plus fréquents et des premiers topics obtenus. Elle vise à améliorer la qualité interprétative de la matrice terme-document en conservant prioritairement les mots porteurs d'information thématique.

\begin{lstlisting}[language=Python, caption={Fonction de lemmatisation avec exclusion de mots spécifiques}]

def nettoyer_texte(texte):
    doc = nlp(texte)
    tokens = []

    for token in doc:
        lemme = token.lemma_.lower()

        if (
            not token.is_stop
            and not token.is_punct
            and not token.like_num
            and not token.is_space
            and token.pos_ in {"NOUN", "ADJ"}
            and len(lemme) > 2
            and lemme not in stop_word
        ):
            tokens.append(lemme)

    return " ".join(tokens)
\end{lstlisting}

Cette fonction applique le modèle \texttt{spaCy} à chaque segment du corpus afin d'en extraire une version lemmatisée et filtrée. Pour chaque token, le lemme est d'abord converti en minuscules afin de normaliser les formes lexicales. La fonction exclut ensuite les mots vides, la ponctuation, les nombres, les espaces, ainsi que les lemmes trop courts ou présents dans la liste personnalisée de mots à retirer. Le filtrage conserve uniquement les noms et les adjectifs, car ces catégories grammaticales sont généralement les plus porteuses d'information thématique dans le cadre d'une analyse de topics. Le résultat final est une chaîne de caractères composée des lemmes retenus, qui peut ensuite être utilisée pour construire la matrice terme-document.

\newpage
\section{Analyse avec NMF et K-means}

Comme mentionné précédemment, deux méthodes d'analyse ont été utilisées pour identifier les thèmes dans les romans de Zola : le NMF (Non-negative Matrix Factorization) et K-means. Je vais expliquer en détail ces méthodes et présenter les résultats obtenus à partir de l'application de ces techniques sur le corpus préparé.

\subsection{La méthode NMF}

Le NMF est une technique de réduction de dimensionnalité qui permet d'extraire des thèmes à partir d'un corpus de textes. Elle fonctionne en factorisant la matrice terme-document en deux matrices non négatives : une matrice de thèmes et une matrice de poids. La matrice de thèmes contient les mots les plus représentatifs de chaque thème, tandis que la matrice de poids indique l'importance de chaque thème dans chaque document (ou paquet de phrases).

C'est la premiere des dux méthodes que j'ai utilisé pour identifier les thèmes dans les romans de Zola.

Après la préparation du corpus, la segmentation en paquets de phrases et la lemmatisation des textes, l'étape suivante consiste à appliquer une méthode de modélisation thématique. L'objectif est d'identifier automatiquement les grands thèmes présents dans les romans de Zola, à partir des termes les plus représentatifs de chaque segment textuel.

Dans cette analyse, la méthode retenue est la NMF, pour \textit{Non-negative Matrix Factorization}. Cette méthode permet de décomposer une matrice terme-document en deux matrices plus simples : une matrice associant les segments de texte aux topics, et une matrice associant les topics aux termes les plus caractéristiques. La NMF est particulièrement adaptée à ce type d'analyse, car elle produit des résultats relativement interprétables : chaque topic peut être compris comme un ensemble de mots fortement associés entre eux.

\subsubsection{Création de la matrice terme-document}

La première étape consiste à transformer les textes lemmatisés en une représentation numérique exploitable par le modèle. Pour cela, une matrice terme-document a été construite à l'aide d'une pondération TF-IDF\footnote{Utilisé via la biblothèque de scikit learn}. Cette méthode permet de donner un poids plus important aux termes caractéristiques d'un segment, tout en réduisant l'influence des mots trop fréquents dans l'ensemble du corpus.

Le vectoriseur utilisé est défini de la manière suivante :

\begin{lstlisting}[language=Python, caption={Création de la matrice terme-document avec TF-IDF}]
vectorizer = TfidfVectorizer(
    min_df=3,
    max_df=0.8
)
X = vectorizer.fit_transform(df["phrases_lemm"])
\end{lstlisting}

Les paramètres \texttt{min\_df} et \texttt{max\_df} jouent un rôle important dans la qualité de la matrice obtenue. Le paramètre \texttt{min\_df=3} permet d'exclure les mots qui apparaissent dans moins de trois segments. Ces mots sont généralement trop rares pour contribuer efficacement à l'identification de thèmes récurrents. À l'inverse, le paramètre \texttt{max\_df=0.8} exclut les termes présents dans plus de 80 \% des segments, car ils sont trop fréquents pour être réellement discriminants.

Ce choix relève d'un arbitrage méthodologique. Plusieurs combinaison ont été testées, puis évaluées à partir de la qualité des mots conservés et de la cohérence des premiers topics obtenus. L'objectif n'était donc pas seulement de réduire la taille de la matrice, mais surtout de conserver les termes les plus utiles pour l'analyse thématique. La matrice finale sert ensuite de base à l'application du modèle NMF.

\subsubsection{Choix du nombre de topics}

Le choix du nombre de topics constitue une étape importante dans une analyse par NMF. Un nombre trop faible de topics risque de produire des catégories trop générales, qui regroupent plusieurs thèmes différents dans un même ensemble. À l'inverse, un nombre trop élevé de topics peut conduire à une fragmentation excessive des résultats, avec des thèmes trop proches les uns des autres ou difficiles à interpréter.

Afin d'orienter ce choix, une méthode du coude a d'abord été utilisée. Elle consiste à entraîner plusieurs modèles NMF avec différents nombres de topics, puis à observer l'évolution de l'erreur de reconstruction. L'idée est de repérer un point à partir duquel l'ajout de nouveaux topics améliore moins fortement la qualité de reconstruction.

\begin{lstlisting}[language=Python, caption={Recherche du nombre de topics par la méthode du coude}]
errors = []
k_values = range(20, 40)

for k in k_values:
    nmf = NMF(
        n_components=k,
        random_state=42,
        init="nndsvda",
        max_iter=500
    )

    nmf.fit(X)
    errors.append(nmf.reconstruction_err_)
\end{lstlisting}

Dans le cas de ce corpus, la méthode du coude ne permet pas d'identifier un seuil clair. L'erreur de reconstruction diminue progressivement lorsque le nombre de topics augmente, mais sans rupture nette. Cela signifie qu'il n'existe pas de valeur optimale évidente à partir de ce seul critère quantitatif.

Le choix du nombre de topics a donc été complété par une analyse qualitative des résultats. Plusieurs valeurs ont été testées, notamment 10, 15, 20, 30 et 35 topics. Les résultats obtenus ont ensuite été comparés en observant les mots caractéristiques de chaque topic, leur cohérence interne et leur capacité à correspondre à des thèmes interprétables dans l'œuvre de Zola.

\subsubsection{Paramétrage du modèle NMF}

Après plusieurs essais, le modèle final a été construit avec 29 topics. Cette valeur a été retenue car elle permettait d'obtenir un niveau de détail suffisant, sans produire une fragmentation trop importante des thèmes (le bruit). Les topics obtenus correspondent à des ensembles thématiques distincts, comme le monde ouvrier, la religion, la famille, la finance, la guerre, le commerce, la médecine ou encore l'espace domestique comme nous allons le voir dans la section d'interprétation des résultats.

Le modèle final est défini de la manière suivante :

\begin{lstlisting}[language=Python, caption={Entraînement du modèle NMF final}]
nmf = NMF(
    n_components=29,
    random_state=42,
    init="nndsvda",
    solver="cd",
    max_iter=1000
)

W = nmf.fit_transform(X)
H = nmf.components_
\end{lstlisting}

Le paramètre \texttt{n\_components=29} correspond au nombre de topics recherchés. Le paramètre \texttt{random\_state=42} permet d'assurer la reproductibilité des résultats. L'initialisation \texttt{nndsvda}\footnote{NNDSVD with zeros filled with the average of X (better when sparsity is not desired) https://scikit-learn.org/stable/modules/generated/sklearn.decomposition.NMF.html } est adaptée à la NMF, car elle fournit une initialisation stable à partir de la structure de la matrice. Le solveur \texttt{cd}, pour \textit{coordinate descent}, est utilisé pour optimiser la factorisation.

La matrice \texttt{W} contient les poids des topics pour chaque segment du corpus. Autrement dit, elle indique dans quelle mesure chaque segment est associé à chacun des 29 topics. La matrice \texttt{H}, quant à elle, contient les poids des mots pour chaque topic. Elle permet donc d'identifier les termes les plus représentatifs de chaque thème.

\subsubsection{Identification et interprétation des topics}

Une fois le modèle entraîné, les mots les plus importants de chaque topic ont été extraits afin de faciliter leur interprétation. Cette étape est essentielle, car la NMF ne produit pas directement des noms de thèmes : elle produit des groupes de mots. L'attribution d'un nom à chaque topic repose donc sur une analyse qualitative des termes les plus représentatifs.

\begin{lstlisting}[language=Python, caption={Affichage des mots les plus importants par topic}]
def afficher_topics(model, feature_names, n_top_words=10):
    for topic_idx, topic in enumerate(model.components_):
        top_words = [
            feature_names[i]
            for i in topic.argsort()[:-n_top_words - 1:-1]
        ]
        print(f"Topic {topic_idx+1} : {' | '.join(top_words)}")

feature_names = vectorizer.get_feature_names_out()
afficher_topics(nmf, feature_names)
\end{lstlisting}

Chaque topic a ensuite été nommé à partir des mots dominants et de l'observation des segments les plus fortement associés à ce topic. Cette double vérification permet de limiter les erreurs d'interprétation. En effet, certains mots peuvent sembler cohérents isolément, mais prendre un sens différent lorsqu'ils sont replacés dans les extraits du corpus.

Les topics obtenus couvrent plusieurs dimensions importantes de l'œuvre de Zola. Certains renvoient à des espaces sociaux précis, comme le magasin, la mine, la bourse, le quartier ou le monde ferroviaire. D'autres correspondent à des thèmes plus transversaux, comme la famille, le désir, la maladie, la religion, la pauvreté ou la mort. Cette diversité montre que la NMF permet de faire apparaître à la fois des thèmes sociaux, des motifs narratifs et des univers propres à certains romans.

\subsubsection{Attribution d'un topic dominant à chaque segment}

Pour analyser plus précisément la répartition des thèmes dans le corpus, un topic dominant a été attribué à chaque segment. Pour cela, les poids de la matrice \texttt{W} ont d'abord été normalisés afin de rendre les contributions des topics comparables entre les segments. Le topic dominant correspond ensuite au topic ayant le poids le plus élevé pour un segment donné.

\begin{lstlisting}[language=Python, caption={Attribution du topic dominant à chaque segment}]
W_norm = W / W.sum(axis=1, keepdims=True)

df["topic_dominant"] = W_norm.argmax(axis=1)
df["poids_topic"] = W_norm.max(axis=1)
\end{lstlisting}

La colonne \texttt{topic\_dominant} indique le topic le plus représenté dans chaque paquet de phrases. La colonne \texttt{poids\_topic} mesure l'intensité de cette association. Plus cette valeur est élevée, plus le segment est clairement associé à un topic spécifique. À l'inverse, un poids faible peut indiquer un segment plus mixte, contenant plusieurs thèmes ou ne correspondant fortement à aucun topic.

Cette étape permet de passer d'une simple liste de topics à une analyse plus structurée de leur répartition dans le corpus. Elle rend possible l'étude du nombre de segments associés à chaque topic, du poids moyen des topics dans l'ensemble du corpus et de la présence des thèmes selon les romans. Et par la suite de passer un detecteur de mofifs sur les segments associés à chaque topic afin d'identifier les motifs narratifs les plus caractéristiques de chaque thème.

\subsubsection{Nommer les topics}

Les topics ont ensuite été nommés manuellement à partir des mots les plus représentatifs et des segments les plus fortement associés. Cette étape constitue une part qualitative importante de l'analyse. Elle permet de transformer les résultats numériques du modèle en catégories interprétables.

\begin{lstlisting}[language=Python, caption={Attribution manuelle des noms de topics}]
noms_topics_nmf = {
    1: "Les relations conjugales et la vie domestique",
    2: "Le travail, le peuple et la justice sociale",
    3: "La nature et le paysage",
    4: "La bourse, la finance et les affaires",
    5: "Le clerge et la paroisse",
    6: "La table et les repas",
    7: "Le corps, la souffrance et la mort",
    8: "L'argent et la fortune",
    9: "Les conversations et les scenes domestiques",
    10: "Le quartier et la ville",
    11: "Le monde ferroviaire",
    12: "Le miracle, le pelerinage et la religion",
    13: "Les liens familiaux, l'ecriture et les proches",
    14: "Le pouvoir religieux",
    15: "L'ecole et l'instruction",
    16: "La famille, la terre et le monde paysan",
    17: "Le monde ouvrier et la mine",
    18: "La guerre, les soldats et l'armee",
    19: "Le desir, le mariage et les relations amoureuses",
    20: "Le magasin, la vente et les clientes",
    21: "Le monde medical et la sante",
    22: "Le monde artistique et la peinture",
    23: "La chambre, le lit et l'espace intime",
    24: "L'enfance, la maternite et la pauvrete",
    25: "Le commerce alimentaire et les halles",
    26: "Le salon, l'aristocratie et les scenes mondaines",
    27: "Le pouvoir politique et les affaires publiques",
    28: "L'amour, le bonheur et les sentiments",
    29: "L'insurrection et la revolution"
}
\end{lstlisting}



Cette étape montre que l’analyse ne repose pas uniquement sur une extraction automatique. Le modèle propose des regroupements de mots, mais leur interprétation dépend d’un travail humain. Les noms attribués aux topics doivent donc être compris comme des catégories d’analyse construites à partir des résultats du modèle, et non comme des vérités objectives produites automatiquement par l’algorithme.


\subsubsection{Limites de l’approche}

L’utilisation de la NMF présente plusieurs avantages, notamment sa lisibilité et sa capacité à produire des topics relativement interprétables. Cependant, cette méthode comporte aussi des limites. Le choix du nombre de topics reste en partie subjectif, surtout lorsque la méthode du coude ne permet pas d’identifier une valeur optimale claire. Dans cette analyse, le choix de 29 topics repose donc sur un compromis entre précision, lisibilité et cohérence interprétative. La validation humaine des resultats est donc essentielle pour s'assurer que les topics produits correspondent à des thèmes réels et pertinents dans l'œuvre de Zola.

De plus, les noms attribués aux topics résultent encore une fois d’une interprétation humaine. Deux personnes pourraient parfois proposer des intitulés légèrement différents pour un même groupe de mots. Il est donc important de considérer ces noms comme des outils d’analyse, et non comme des catégories fixes ou définitives. Cette subjectivité ne remet pas en cause la validité de l’analyse, mais elle souligne l’importance d’une approche critique et réflexive dans l’interprétation des résultats.

Enfin, les résultats dépendent fortement des choix effectués en amont : nettoyage du texte, lemmatisation, liste de mots à exclure, taille des segments, paramètres du vectoriseur et nombre de topics. La modélisation thématique ne doit donc pas être comprise comme une méthode entièrement automatique, mais comme un processus exploratoire qui combine traitement statistique et interprétation qualitative avec validation humaine à chaque étapes.

\subsubsection{Bilan méthodologique}

Cette étape de modélisation par NMF permet d’obtenir une première cartographie thématique du corpus romanesque de Zola. Les résultats montrent que les romans peuvent être décrits à partir de grands ensembles thématiques cohérents, liés à la fois aux espaces sociaux, aux relations familiales, aux institutions, aux métiers, aux croyances et aux conflits collectifs.

La NMF constitue donc un outil pertinent pour explorer les thèmes récurrents dans un corpus littéraire volumineux. Elle ne remplace pas l’analyse littéraire, mais elle permet de faire apparaître des régularités lexicales et thématiques qui peuvent ensuite être interprétées plus finement. Dans le cadre de cette étude, elle sert à établir une base quantitative pour l’analyse des thèmes présents dans les romans de Zola.

\subsection{La méthode du K-means}

Après l'application de la NMF, une seconde méthode a été utilisée afin de compléter l'analyse thématique du corpus : le clustering par K-means. Contrairement à la NMF, qui cherche à décomposer la matrice terme-document en topics latents, K-means vise à regrouper les segments textuels en ensembles homogènes à partir de leur proximité lexicale. L'objectif n'est donc pas exactement le même : la NMF permet d'identifier des thèmes sous forme de combinaisons de mots, tandis que K-means permet de regrouper les paquets de phrases qui présentent des profils lexicaux similaires.

Cette seconde méthode a été utilisée comme un outil de comparaison et de validation. Elle permet de vérifier si les regroupements produits par la NMF correspondent également à des ensembles cohérents lorsque l'on applique une méthode de clustering non supervisée. Dans le cadre de cette analyse, chaque paquet de phrases constitue une unité d'observation, représentée numériquement par une matrice TF-IDF. Les clusters obtenus correspondent donc à des groupes de segments partageant des termes caractéristiques.

\subsubsection{Préparation de la matrice TF-IDF}

Comme pour la NMF, la première étape consiste à transformer les textes lemmatisés en une représentation numérique. La matrice utilisée pour le K-means est construite à partir des segments lemmatisés du corpus. Le même type de vectorisation TF-IDF et les meme parametres ont été conservé afin d'assurer une cohérence méthodologique entre les deux approches\footnote{Cf: "preparation de la matrice NMF" dans l'approche NMF}.

\begin{lstlisting}[language=Python, caption={Création de la matrice TF-IDF pour le clustering K-means}]
vectorizer = TfidfVectorizer(
    min_df=3,
    max_df=0.8
)
X = vectorizer.fit_transform(df["phrases_lemm"])
\end{lstlisting}

Ce choix permet de travailler sur une matrice plus lisible et plus pertinente pour le clustering. Le but n'est pas seulement de réduire le nombre de dimensions, mais aussi de conserver les mots qui participent réellement à la différenciation des segments. Cette étape est essentielle, car la qualité des clusters dépend directement de la qualité de la représentation vectorielle utilisée.

\subsubsection{Recherche du nombre de clusters}

Le choix du nombre de clusters est une étape centrale dans l'utilisation de K-means. Un nombre trop faible de clusters risque de produire des groupes trop généraux, dans lesquels plusieurs thèmes différents seraient mélangés. À l'inverse, un nombre trop élevé peut produire des groupes trop fragmentés, difficiles à interpréter ou très proches les uns des autres.

Afin d'orienter ce choix, plusieurs valeurs de \texttt{k} ont été testées. Pour chaque valeur comprise entre 2 et 40, un modèle K-means a été entraîné, puis évalué à l'aide du score de silhouette\footnote{Peter J. Rousseeuw, « Silhouettes: A graphical aid to the interpretation and validation of cluster analysis », \textit{Journal of Computational and Applied Mathematics}, vol. 20, 1er novembre 1987, p. 53–65 }. 

Ce score mesure la qualité de la séparation entre les clusters : plus il est élevé, plus les segments sont proches des éléments de leur propre cluster et éloignés des autres clusters.

\begin{lstlisting}[language=Python, caption={Évaluation du nombre de clusters par score de silhouette}]
scores = []

for k in range(2, 41):
    model = KMeans(
        n_clusters=k,
        random_state=42,
        n_init=10
    )
    labels = model.fit_predict(X)
    score = silhouette_score(X, labels)
    scores.append((k, score))

scores_df = pd.DataFrame(scores, columns=["k", "silhouette"])
\end{lstlisting}

Une visualisation des scores de silhouette a ensuite été produite afin d'observer l'évolution de la qualité du clustering selon le nombre de clusters.

\begin{lstlisting}[language=Python, caption={Visualisation du score de silhouette selon le nombre de clusters}]
plt.figure(figsize=(10, 6))
plt.xticks(range(2, 41))

plt.plot(scores_df["k"], scores_df["silhouette"], marker="o")

plt.xlabel("Nombre de clusters k")
plt.ylabel("Score de silhouette")
plt.title("Score de silhouette selon le nombre de clusters")

plt.show()
\end{lstlisting}

Dans cette analyse, le score de silhouette fournit une indication quantitative, mais il ne suffit pas à lui seul à déterminer le nombre de clusters à retenir. Comme pour la NMF, le choix final repose donc sur un compromis entre indicateurs statistiques, lisibilité des résultats et cohérence interprétative. Le nombre de 29 clusters a été retenu afin de permettre une comparaison directe avec les 29 topics issus de la NMF, tout en conservant un niveau de détail suffisant pour distinguer les principaux motifs thématiques du corpus, de plus le score de silhouette reste relativement stable autour de cette valeur, ce qui indique que les clusters obtenus sont globalement cohérents comme pour la NMF.

\subsubsection{Entraînement du modèle K-means}

Le modèle K-means final a été entraîné avec 29 clusters. Le paramètre \texttt{random\_state=42} permet d'assurer la reproductibilité des résultats, tandis que le paramètre \texttt{n\_init=10} indique que l'algorithme est lancé plusieurs fois avec des initialisations différentes afin de retenir la meilleure solution.

\begin{lstlisting}[language=Python, caption={Entraînement du modèle K-means final}]
kmeans = KMeans(
    n_clusters=29,
    random_state=42,
    n_init=10
)

df["cluster"] = kmeans.fit_predict(X) + 1
\end{lstlisting}

Chaque segment du corpus reçoit ainsi un numéro de cluster. L'ajout de \texttt{+ 1} permet de numéroter les clusters à partir de 1 plutôt qu'à partir de 0, ce qui rend les résultats plus lisibles dans les tableaux et les visualisations. Chaque paquet de phrases est donc associé à un groupe thématique déterminé par sa proximité avec les autres segments du corpus.

Contrairement à la NMF, qui attribue à chaque segment un poids pour chacun des topics, K-means attribue chaque segment à un seul cluster. Cette différence est importante : la NMF autorise une lecture plus graduelle des thèmes, tandis que K-means produit une classification plus nette. Un segment appartient à un cluster principal, même s'il peut contenir plusieurs motifs thématiques.

\subsubsection{Extraction des mots caractéristiques des clusters}

Une fois le modèle entraîné, les centres des clusters ont été analysés afin d'identifier les termes les plus représentatifs de chaque groupe. Dans K-means, chaque cluster est représenté par un centroïde, c'est-à-dire un point moyen dans l'espace vectoriel. Les termes les plus fortement associés à ce centroïde permettent de comprendre le contenu thématique du cluster.

\begin{lstlisting}[language=Python, caption={Extraction des mots les plus représentatifs par cluster}]
terms = vectorizer.get_feature_names_out()
centres = kmeans.cluster_centers_

for i, centre in enumerate(centres):
    top_indices = centre.argsort()[::-1][:10]
    mots = [terms[j] for j in top_indices]
    print(f"Cluster {i+1} :", " | ".join(mots))
\end{lstlisting}

Cette étape est comparable à l'extraction des mots dominants dans la NMF. Elle permet d'associer chaque cluster à un ensemble de termes caractéristiques. Cependant, l'interprétation est légèrement différente : dans la NMF, les mots sont associés à un topic latent ; dans K-means, ils décrivent le centre lexical d'un groupe de segments. Autrement dit, les mots obtenus indiquent ce qui rapproche les segments d'un même cluster.

Cette extraction permet de vérifier la cohérence interne des clusters. Lorsqu'un cluster rassemble des mots appartenant au même champ lexical, il devient plus facilement interprétable. À l'inverse, un cluster contenant des mots trop dispersés ou peu cohérents peut indiquer un regroupement moins stable ou plus difficile à nommer.

\subsubsection{Nomination des clusters}

Les clusters ont ensuite été nommés manuellement à partir des mots les plus représentatifs et de l'observation des segments associés. Cette étape est indispensable, car K-means ne produit pas directement des intitulés interprétables. Comme pour la NMF, le passage du résultat numérique à l'interprétation thématique repose donc sur une lecture qualitative.

\begin{lstlisting}[language=Python, caption={Attribution manuelle des noms de clusters}]
noms_clusters_kmeans = {
    1: "La table, les repas et les boissons",
    2: "L'ecole, les eleves et l'instruction",
    3: "Le salon, l'aristocratie et les apparences",
    4: "La chambre, le lit et l'espace domestique",
    5: "L'enfance, la maternite et la pauvrete",
    6: "La mere, les enfants et la famille",
    7: "Le monde ferroviaire",
    8: "L'amour, les sentiments et la jeunesse",
    9: "Le theatre, la loge et les scenes mondaines",
    10: "Le clerge et la paroisse",
    11: "La nature et le paysage",
    12: "Le miracle, le pelerinage et la religion",
    13: "Le monde medical et la sante",
    14: "Le pelerinage, la maladie et le voyage",
    15: "L'insurrection et la revolution",
    16: "Le quartier, la rue et la ville",
    17: "Le desir, le corps et les relations amoureuses",
    18: "Le pouvoir politique et les affaires publiques",
    19: "Le travail, le peuple et le monde ouvrier",
    20: "La bourse, la finance et les affaires",
    21: "L'argent du quotidien",
    22: "Les conversations et les scenes domestiques",
    23: "Le monde artistique et la peinture",
    24: "La mine, les ouvriers et la greve",
    25: "Le magasin, la vente et les clientes",
    26: "La guerre, les soldats et l'armee",
    27: "Le pouvoir religieux",
    28: "Les liens familiaux, l'ecriture et la justice",
    29: "La terre, les animaux et le monde paysan"
}

df["nom_cluster"] = df["cluster"].map(noms_clusters_kmeans)
\end{lstlisting}

Les intitulés obtenus montrent que les clusters couvrent des dimensions proches de celles observées avec la NMF. On retrouve par exemple des ensembles liés au travail, à la mine, à la finance, à la religion, à la famille, au monde domestique, à la guerre ou encore aux espaces urbains. Cette convergence entre les deux méthodes est intéressante, car elle indique que certains thèmes apparaissent de manière stable dans le corpus, indépendamment de la méthode utilisée.

Cependant, les clusters K-means ne doivent pas être interprétés comme des catégories fixes ou naturelles. Ils résultent d'un regroupement algorithmique fondé sur la proximité lexicale des segments. Leur nomination repose donc sur un choix interprétatif, à partir des mots les plus représentatifs et des extraits textuels associés.

\subsubsection{Regroupement hiérarchique des clusters}

Afin de mieux comprendre les relations entre les 29 clusters, une classification ascendante hiérarchique a été appliquée aux centres des clusters. Cette étape permet de regrouper les clusters proches entre eux et d'identifier des familles thématiques plus larges. Elle complète donc le K-means en donnant une lecture plus synthétique de la structure thématique du corpus.

\begin{lstlisting}[language=Python, caption={Classification hiérarchique sur les centres des clusters}]
Z = linkage(centres, method="ward")

dendrogram(
    Z,
    labels=[f"Cluster {i+1}" for i in range(kmeans.n_clusters)],
    orientation="right",
    color_threshold=0.40
)

plt.axvline(x=0.40, color="red", linestyle="--")

plt.title("CAH sur les clusters K-means")
plt.xlabel("Distance")
plt.show()
\end{lstlisting}

La méthode de Ward a été utilisée pour construire la hiérarchie. Elle cherche à regrouper les clusters en minimisant l'augmentation de la variance interne à chaque fusion. Le dendrogramme permet ensuite de visualiser les proximités entre clusters. Le seuil de distance fixé à 0,40 sert à découper l'arbre en grandes familles de motifs.

Cette étape est utile car 29 clusters peuvent être difficiles à interpréter globalement. Le regroupement hiérarchique permet de passer d'une lecture détaillée, cluster par cluster, à une lecture plus générale, organisée autour de grandes familles thématiques.

\subsubsection{Création des grandes familles de motifs}

À partir du dendrogramme, les clusters ont été regroupés en cinq grandes familles de motifs. Cette opération permet de simplifier l'analyse et de faire apparaître les grandes zones thématiques du corpus.

\begin{lstlisting}[language=Python, caption={Création des groupes CAH à partir des clusters K-means}]
groupes_cah = fcluster(Z, t=0.40, criterion="distance")

df_groupes = pd.DataFrame({
    "cluster": range(1, kmeans.n_clusters + 1),
    "groupe_cah": groupes_cah
})

df_groupes["nom_cluster"] = df_groupes["cluster"].map(
    noms_clusters_kmeans
)
\end{lstlisting}

Les groupes obtenus ont ensuite été nommés manuellement, en fonction des clusters qu'ils rassemblent.

\begin{lstlisting}[language=Python, caption={Nomination des grandes familles de motifs}]
noms_groupes_cah = {
    1: "Religion locale et autorite clericale",
    2: "Guerre, insurrection et conflit politique",
    3: "Institutions, croyances et encadrement social",
    4: "Argent, dettes et circulation financiere",
    5: "Vie quotidienne, espaces sociaux et experiences individuelles"
}

df_groupes["nom_groupe_cah"] = df_groupes["groupe_cah"].map(
    noms_groupes_cah
)
\end{lstlisting}

Ces grandes familles permettent d'obtenir une lecture plus synthétique des résultats. La première famille regroupe les thèmes liés à la religion locale et à l'autorité cléricale. La deuxième rassemble les motifs de guerre, d'insurrection et de conflit politique. La troisième renvoie aux institutions, aux croyances et aux formes d'encadrement social. La quatrième concerne l'argent, les dettes et les circulations financières. Enfin, la cinquième regroupe les motifs les plus liés à la vie quotidienne, aux espaces sociaux et aux expériences individuelles.

Cette structuration est particulièrement utile pour l'analyse littéraire, car elle permet de passer d'une liste de clusters nombreux à une organisation thématique plus lisible. Elle montre aussi que les motifs extraits par K-means ne sont pas isolés, mais peuvent être regroupés en ensembles plus larges correspondant à des dimensions importantes de l'œuvre de Zola.

\subsubsection{Attribution des grandes familles aux segments}

Une fois les groupes hiérarchiques définis, chaque segment du corpus a reçu une étiquette correspondant à la grande famille de motifs associée à son cluster. Cette opération permet d'analyser non seulement les clusters détaillés, mais aussi les grandes catégories thématiques dans lesquelles ils s'inscrivent.

\begin{lstlisting}[language=Python, caption={Attribution des groupes CAH aux segments du corpus}]
cluster_vers_groupe = dict(
    zip(df_groupes["cluster"], df_groupes["groupe_cah"])
)

df["groupe_cah"] = df["cluster"].map(cluster_vers_groupe)
df["nom_groupe_cah"] = df["groupe_cah"].map(noms_groupes_cah)
\end{lstlisting}

Cette étape permet de relier chaque paquet de phrases à deux niveaux d'analyse. Le premier niveau est le cluster K-means, qui fournit une catégorie thématique fine. Le second niveau est le groupe CAH, qui fournit une catégorie plus générale. Cette double échelle d'interprétation est utile car elle permet d'adapter le niveau de détail selon les besoins de l'analyse : on peut étudier les motifs précis, ou bien observer les grandes tendances du corpus.

\subsubsection{Analyse de la répartition des clusters}

Après la nomination des clusters et des grandes familles, plusieurs visualisations ont été produites afin d'étudier leur répartition dans le corpus. La première consiste à compter le nombre de segments associés à chaque grande famille de motifs.

\begin{lstlisting}[language=Python, caption={Nombre de segments par grande famille de motifs}]
counts = df["nom_groupe_cah"].value_counts().sort_values()

plt.figure(figsize=(10, 6))
counts.plot(kind="barh")

plt.title("Nombre de segments par grande famille de motifs")
plt.xlabel("Nombre de segments")
plt.ylabel("Groupe CAH")
plt.show()
\end{lstlisting}

Cette visualisation permet d'identifier les grandes familles les plus représentées dans l'ensemble du corpus. Elle donne une lecture globale de l'organisation thématique des segments. Une famille très représentée indique qu'un ensemble de motifs revient fréquemment dans les romans, tandis qu'une famille moins présente peut correspondre à un thème plus spécifique, associé à certains romans ou à certaines situations narratives.

Une seconde visualisation a été produite à l'échelle des clusters détaillés.


Cette représentation permet de voir quels clusters sont les plus fréquents dans le corpus. Elle est plus précise que l'analyse par grandes familles, car elle distingue les motifs thématiques particuliers. Elle permet par exemple d'observer séparément les clusters associés à la mine, à la finance, à la religion, à la famille ou au monde domestique.

\subsubsection{Distribution des clusters par œuvre}

Afin d'observer la présence des clusters dans chaque roman, une table croisée a été construite entre les fichiers du corpus et les clusters thématiques. La normalisation par ligne permet d'obtenir, pour chaque œuvre, la proportion de segments associés à chaque cluster.



Cette analyse est importante car elle permet de relier les clusters aux romans. Certains clusters peuvent être très transversaux et apparaître dans de nombreuses œuvres, tandis que d'autres peuvent être fortement associés à un roman précis. Par exemple, un cluster lié à la mine peut être particulièrement représenté dans \textit{Germinal}, tandis qu'un cluster lié à la finance peut être plus présent dans \textit{L'Argent}. De la même manière, les clusters associés au commerce peuvent être importants dans \textit{Au Bonheur des dames}, et ceux liés à la guerre dans \textit{La Débâcle}.

Une analyse similaire a été réalisée au niveau des grandes familles de motifs.



Cette seconde visualisation offre une lecture plus synthétique. Elle permet de comparer les romans selon les grandes familles de motifs qui les structurent. Cette approche est particulièrement utile pour repérer des oppositions ou des proximités entre les œuvres : certains romans peuvent être dominés par des motifs de vie quotidienne, tandis que d'autres accordent plus de place aux institutions, à la religion, à la guerre ou aux conflits sociaux.

\subsubsection{Projection UMAP des segments}

Une projection UMAP a enfin été utilisée pour représenter les segments dans un espace en deux dimensions. Cette méthode de réduction de dimensionnalité permet de visualiser la manière dont les segments se répartissent selon leur proximité lexicale. Les points sont colorés selon leur grande famille de motifs.

\begin{lstlisting}[language=Python, caption={Projection UMAP des segments par grande famille de motifs}]
reducer = umap.UMAP(n_components=2, random_state=42)
X_umap = reducer.fit_transform(X)

plt.figure(figsize=(10, 7))

plt.scatter(
    X_umap[:, 0],
    X_umap[:, 1],
    c=df["groupe_cah"],
    cmap="tab10",
    s=10,
    alpha=0.7
)

plt.title("Projection UMAP des segments par grande famille de motifs")
plt.xlabel("UMAP 1")
plt.ylabel("UMAP 2")
plt.colorbar(label="Groupe CAH")
plt.show()
\end{lstlisting}

Cette visualisation ne doit pas être interprétée comme une preuve définitive de séparation entre les groupes, mais comme un outil exploratoire. Elle permet d'observer si les segments associés à une même famille de motifs tendent à se rapprocher dans l'espace vectoriel. Si certaines zones apparaissent relativement regroupées, cela peut indiquer que les familles thématiques correspondent à des structures lexicales cohérentes. À l'inverse, des zones plus mélangées peuvent signaler des motifs plus transversaux ou des segments hybrides.

\subsubsection{Limites de l'approche par K-means}

L'utilisation de K-means présente plusieurs limites qu'il faut prendre en compte. Tout d'abord, la méthode impose une appartenance unique de chaque segment à un cluster. Or, dans un texte littéraire, un même passage peut contenir plusieurs thèmes simultanément. Un segment peut par exemple associer la famille, l'argent et la souffrance sociale, mais K-means l'attribue malgré tout à un seul groupe. Cette contrainte rend la méthode moins souple que la NMF pour représenter la complexité thématique des textes.

Ensuite, le choix du nombre de clusters reste en partie subjectif. Le score de silhouette fournit une indication utile, mais il ne suffit pas toujours à identifier une valeur optimale. Dans cette analyse, le choix de 29 clusters repose donc sur un arbitrage méthodologique : il permet de conserver un niveau de détail comparable à celui de la NMF, tout en produisant des catégories suffisamment interprétables.

Une autre limite concerne la sensibilité de K-means à la représentation vectorielle. Les clusters dépendent fortement de la matrice TF-IDF, des choix de lemmatisation, de la liste de mots exclus et des paramètres \texttt{min\_df} et \texttt{max\_df}. Une modification de ces choix peut modifier la structure des clusters obtenus tout comme la méthode du NMF. Les résultats doivent donc être compris comme une cartographie exploratoire du corpus, et non comme une classification définitive.

Enfin, comme pour la NMF, la nomination des clusters repose sur une interprétation humaine. Les intitulés attribués aux clusters et aux grandes familles de motifs sont construits à partir des mots les plus représentatifs et de l'observation des segments. Cette étape qualitative est nécessaire pour donner du sens aux résultats, mais elle introduit aussi une part de subjectivité.

\subsubsection{Bilan méthodologique}

La méthode K-means permet de compléter l'analyse réalisée avec la NMF. Elle offre une autre manière d'observer les régularités thématiques du corpus, non plus sous la forme de topics latents, mais sous la forme de groupes de segments lexicalement proches. Cette approche confirme plusieurs motifs déjà observés avec la NMF, notamment autour du travail, de la religion, de la famille, de l'argent, de la guerre, de la mine, du commerce ou de l'espace domestique.

L'ajout d'une classification hiérarchique sur les centres des clusters permet également d'organiser les résultats en grandes familles de motifs. Cette étape rend l'analyse plus lisible, car elle permet de passer de 29 clusters détaillés à cinq ensembles thématiques plus généraux. Ces familles offrent une vision synthétique de la structure thématique du corpus et facilitent la comparaison entre les œuvres.

Ainsi, K-means ne remplace pas la NMF, mais il constitue un complément méthodologique utile. En croisant les résultats des deux méthodes, il devient possible de distinguer les thèmes particulièrement stables de ceux qui dépendent davantage du modèle utilisé. Cette double approche renforce donc la solidité de l'analyse, tout en conservant une part d'interprétation qualitative indispensable dans l'étude d'un corpus littéraire.

\end{document}